\documentclass[a4paper,fleqn]{cas-dc}

% ------------------------------------------------------------------
% Citation style
% ------------------------------------------------------------------

\usepackage[authoryear,longnamesfirst]{natbib}

% ------------------------------------------------------------------
% Mathematical and typographical packages
% ------------------------------------------------------------------

\usepackage{amsmath}
\usepackage{amssymb}
\usepackage{amsthm}
\usepackage{mathtools}
\usepackage{booktabs}
\usepackage{graphicx}
\usepackage[
    protrusion=true,
    expansion=false
]{microtype}
\usepackage{xspace}

% ------------------------------------------------------------------
% Theorem environments
% ------------------------------------------------------------------

\newtheorem{proposition}{Proposition}
\newtheorem{corollary}{Corollary}
\newtheorem{remark}{Remark}

% ------------------------------------------------------------------
% Document
% ------------------------------------------------------------------

\begin{document}

\let\WriteBookmarks\relax
\def\floatpagepagefraction{1}
\def\textpagefraction{0.001}

% ------------------------------------------------------------------
% Running title and author
% ------------------------------------------------------------------

\shorttitle{Concentration Is Not Scaling}
\shortauthors{M. Iglesias}

% ------------------------------------------------------------------
% Title
% ------------------------------------------------------------------

\title[mode=title]{
Concentration Is Not Scaling:
An Exact Decomposition of Aggregate Variance Elasticities
}

% ------------------------------------------------------------------
% Author
% ------------------------------------------------------------------

\author[1]{Matías Iglesias}
\cormark[1]

\ead{mniglesias@dc.uba.ar}

\credit{
Conceptualization,
Methodology,
Formal analysis,
Software,
Validation,
Visualization,
Writing -- Original Draft,
Writing -- Review \& Editing
}

% ------------------------------------------------------------------
% Affiliation
% ------------------------------------------------------------------

\affiliation[1]{
    organization={
        Departamento de Computación,
        Facultad de Ciencias Exactas y Naturales,
        Universidad de Buenos Aires
    },
    city={Buenos Aires},
    country={Argentina}
}

\cortext[cor1]{Corresponding author}

% ------------------------------------------------------------------
% Abstract
% ------------------------------------------------------------------

\begin{abstract}
Granular models relate aggregate variance to the concentration of
economic weights. This proportionality is exact only under restrictive
assumptions on marginal shock variances, dependence, and population
composition. Concentration is a level property, whereas variance
scaling is an elasticity defined along a population path. We show that
identical Herfindahl paths can generate different aggregate-variance
elasticities and derive an exact finite-interval decomposition of that
elasticity into changes in diagonal variance and signed covariance
components. The decomposition uses the logarithmic mean of total
variance and requires no logarithm of individual components. A
contribution-weighted average of component elasticities emerges only
under positivity and comparability conditions. Deterministic analytic
constructions verify the direct and reconstructed elasticities at
floating-point precision. The result clarifies the boundary of
concentration-based granular reasoning and supplies an accounting
framework for micro-to-macro analyses of diversification.
\end{abstract}

% ------------------------------------------------------------------
% Highlights
% ------------------------------------------------------------------

\begin{highlights}
\item Concentration does not identify aggregate-variance scaling.
\item Variance elasticities admit an exact finite-interval decomposition.
\item Signed covariance enters without logarithms of component values.
\item Equal concentration paths can imply different diversification rates.
\item Positive covariance can contribute negatively to variance decay.
\end{highlights}

% ------------------------------------------------------------------
% Keywords and classification
% ------------------------------------------------------------------

\begin{keywords}
aggregate variance
\sep concentration
\sep diversification
\sep variance decomposition
\sep covariance
\sep scaling elasticity
\end{keywords}

\nonumnote{
JEL classification: C02, C10, D20.
}

% ------------------------------------------------------------------
% Generate title
% ------------------------------------------------------------------

\maketitle

% ------------------------------------------------------------------
% Main text
% ------------------------------------------------------------------

\section{Introduction}
\label{sec:introduction}

The granular approach restored an elementary but consequential fact to
macroeconomic aggregation: units enter an aggregate according to their
economic weights, so a small number of very large firms can prevent
idiosyncratic fluctuations from averaging out rapidly
\citep{gabaix2011}. Under mutually independent shocks with a common
variance, the variance of a weighted aggregate is proportional to the
Herfindahl index
\begin{equation}
    H(N)=\sum_{i=1}^{N}w_i(N)^2.
\end{equation}
A fat-tailed firm-size distribution can therefore generate a
concentration path that declines slowly as the population expands.

This benchmark is sometimes given a stronger interpretation: that the
concentration path identifies the population-scaling rate of aggregate
variance. That conclusion is not general. Concentration describes the
distribution of weights at a given scale. Aggregate-variance scaling
is an elasticity describing how the complete variance--covariance
structure changes along a specified population path. The two objects
coincide only when restrictions on marginal variances, dependence, and
composition supply the missing bridge between them.

For the weighted aggregate
\begin{equation}
    X_t(N)=\sum_{i=1}^{N}w_i(N)\varepsilon_{i,t}(N),
\end{equation}
aggregate variance is
\begin{align}
\label{eq:intro_variance}
    V(N)
    &=
    \sum_{i=1}^{N}w_i(N)^2
    \operatorname{Var}\!\left[\varepsilon_{i,t}(N)\right]
    \nonumber\\
    &\quad+
    2\sum_{i<j}w_i(N)w_j(N)
    \operatorname{Cov}\!\left[
        \varepsilon_{i,t}(N),\varepsilon_{j,t}(N)
    \right].
\end{align}
The proportionality $V(N)=\sigma^2H(N)$ follows when all shocks share a
variance $\sigma^2$ that is stable along the population path and the
off-diagonal terms vanish. These are not incidental simplifications.
Even under independence, allowing the common shock variance to vary with
population scale gives
\begin{equation}
\label{eq:benchmark_elasticity_bridge}
    -\frac{d\log V(N)}{d\log N}
    =
    -\frac{d\log H(N)}{d\log N}
    -\frac{d\log \sigma^2(N)}{d\log N}.
\end{equation}
The concentration elasticity identifies the variance elasticity only
when the second term is zero. Heterogeneous firm-level variances replace
$H(N)$ by a variance-weighted concentration object, while common shocks
and economic linkages add off-diagonal terms. Population expansion may
also alter the size, sectoral, or network composition of the aggregate.
Once these channels are present, the Herfindahl index is not a sufficient
statistic for either aggregate variance or its population elasticity.

This paper separates the granular insight from that stronger inference.
Granularity correctly establishes that aggregate variance is weighted
by firm shares and may be dominated by a small number of large units. It
does not, by itself, establish that the firm-size distribution identifies
the rate at which aggregate variance changes with population size. We
make this distinction precise in two steps.

First, we show that identical concentration paths can generate different
aggregate-variance scaling. Equal-weight populations all satisfy
$H(N)=1/N$, but aggregate variance follows $1/N$ only under independence.
With persistent positive correlation, it instead approaches a positive
covariance floor and its scaling elasticity converges to zero. No
statistic constructed solely from the weights can distinguish the two
systems.

Second, we derive an exact finite-interval decomposition of the aggregate
variance elasticity into changes in diagonal variance and signed
covariance components. The formula does not take logarithms of individual
components and therefore remains well defined when covariance
contributions are negative. It also identifies the restrictive
positive-component case in which an aggregate elasticity may be read as
a contribution-weighted average of component elasticities.

The mathematical ingredients are established. Granular aggregation
studies unequal economic weights, dependence enters aggregate variance
through covariance, and logarithmic-mean methods provide exact additive
change accounting. Our claim is deliberately narrower than invention of
any of these ingredients. The contribution is their combination for a
distinct micro-to-macro object: the finite-interval elasticity of
aggregate variance with respect to population scale. This formulation
separates a level statement about weighted exposure from a path
statement about diversification, retains signed covariance
contributions, and yields an identity that can be checked directly in
empirical or synthetic data.
% Insert verified granular, network, decomposition, and risk-attribution
% citations during the bibliography pass.

\section{Exact decomposition}
\label{sec:exact_decomposition}

Consider an aggregate composed of $Q$ contribution series,
\begin{equation}
    X_t(N)=\sum_{q=1}^{Q}c_{q,t}(N),
\end{equation}
where $N>0$ indexes population scale. Its variance is
\begin{equation}
\label{eq:variance_accounting}
    V(N)
    =
    \sum_{q=1}^{Q}\operatorname{Var}\!\left[c_{q,t}(N)\right]
    +
    2\sum_{q<r}\operatorname{Cov}\!\left[c_{q,t}(N),c_{r,t}(N)\right].
\end{equation}
Let $\mathcal{K}$ collect the diagonal and doubled
off-diagonal terms,
\begin{align}
    C_{qq}(N)
    &=
    \operatorname{Var}\!\left[c_{q,t}(N)\right], \\
    C_{qr}(N)
    &=
    2\operatorname{Cov}\!\left[c_{q,t}(N),c_{r,t}(N)\right],
    \qquad q<r.
\end{align}
Then
\begin{equation}
\label{eq:component_sum}
    V(N)=\sum_{k\in\mathcal{K}}C_k(N).
\end{equation}
Diagonal components are non-negative, whereas covariance components may
have either sign. The component universe $\mathcal{K}$ is assumed to
retain a comparable meaning at the population levels being contrasted.

For $N_0,N_1>0$, $N_0\neq N_1$, and positive aggregate variance at both
levels, define the finite-interval variance elasticity
\begin{equation}
\label{eq:direct_alpha}
    \widehat{\alpha}(N_0,N_1)
    =
    -\frac{
        \log V(N_1)-\log V(N_0)
    }{
        \log N_1-\log N_0
    }.
\end{equation}
For $a,b>0$, let
\begin{equation}
\label{eq:log_mean}
    L(a,b)
    =
    \begin{cases}
        \displaystyle\frac{a-b}{\log a-\log b}, & a\neq b,\\[1em]
        a, & a=b,
    \end{cases}
\end{equation}
be the logarithmic mean.

\begin{proposition}[Exact finite-interval decomposition]
\label{prop:exact_decomposition}
For each $k\in\mathcal{K}$, define
\begin{equation}
\label{eq:component_alpha_contribution}
    \widehat{\alpha}^{\,C}_k(N_0,N_1)
    =
    -\frac{
        C_k(N_1)-C_k(N_0)
    }{
        L\!\left(V(N_1),V(N_0)\right)
        \left[\log N_1-\log N_0\right]
    }.
\end{equation}
Then
\begin{equation}
\label{eq:exact_alpha_decomposition}
    \widehat{\alpha}(N_0,N_1)
    =
    \sum_{k\in\mathcal{K}}
    \widehat{\alpha}^{\,C}_k(N_0,N_1).
\end{equation}
The identity remains valid when individual covariance components are
negative.
\end{proposition}

\begin{proof}
Equation~\eqref{eq:component_sum} implies
\begin{equation}
    \sum_{k\in\mathcal{K}}
    \left[C_k(N_1)-C_k(N_0)\right]
    =V(N_1)-V(N_0).
\end{equation}
Summing equation~\eqref{eq:component_alpha_contribution} over $k$ and
using
\begin{equation}
    \frac{V(N_1)-V(N_0)}{L(V(N_1),V(N_0))}
    =\log V(N_1)-\log V(N_0)
\end{equation}
produces equation~\eqref{eq:direct_alpha} exactly.
\end{proof}

The decomposition attributes the aggregate elasticity to changes in the
underlying variance and covariance components without requiring their
individual logarithms. Signed covariance terms therefore enter directly,
provided total variance remains positive. The result is exact for each
finite interval and does not require a global power law. A regression
coefficient estimated over many population levels may still provide a
useful summary, but it averages interval-specific elasticities whose
component weights can change across the fitted range. The decomposition
therefore distinguishes an accounting identity from a parametric claim:
it reconstructs what drives the observed change in aggregate variance
without assuming that one invariant exponent governs every scale. In the
differentiable limit,
\begin{equation}
\label{eq:local_signed_decomposition}
    \alpha_{\mathrm{eff}}(N)
    =-
    \frac{d\log V(N)}{d\log N}
    =-
    \sum_{k\in\mathcal{K}}
    \frac{N C_k'(N)}{V(N)}.
\end{equation}
If every component is positive, define
\begin{equation}
    \alpha_k(N)
    =
    -\frac{d\log C_k(N)}{d\log N}.
\end{equation}
Then
\begin{equation}
\label{eq:weighted_local_alpha}
    \alpha_{\mathrm{eff}}(N)
    =
    \sum_{k\in\mathcal{K}}
    \frac{C_k(N)}{V(N)}\alpha_k(N).
\end{equation}
This resembles a weighted average, but its weights generally vary with
$N$. A fixed average of local exponents arises only under additional
restrictions, such as common component elasticities or contribution
shares that remain invariant along the population path. When positive
components decay at different power-law rates, their shares change with
$N$ and the slowest-decaying component eventually dominates. With signed
covariance components the normalized contributions need not be
non-negative, so the representation is not a convex combination and the
aggregate elasticity need not lie within the range of the diagonal
elasticities. Covariance is therefore not merely a correction to the
level of aggregate variance; its own scale dependence can determine the
limiting diversification rate.

Comparability across scales is substantive. If parts appear, disappear,
or change meaning between $N_0$ and $N_1$, the measured elasticity may
combine diversification with composition, partition, or dependence
changes. The identity remains valid only after those changes are
represented through a common accounting structure.

\section{Why concentration does not identify variance scaling}
\label{sec:concentration_scaling}

Consider the weighted aggregate
\begin{equation}
\label{eq:weighted_aggregate}
    X_t(N)
    =
    \sum_{i=1}^{N}w_i(N)\varepsilon_{i,t}(N),
    \qquad
    \sum_{i=1}^{N}w_i(N)=1.
\end{equation}
For equal weights, $w_i(N)=1/N$, the concentration path is
\begin{equation}
\label{eq:equal_weight_concentration}
    H(N)=\frac{1}{N}.
\end{equation}
Suppose first that shocks have unit variance and are mutually
independent. Aggregate variance is then
\begin{equation}
\label{eq:iid_variance}
    V_{\mathrm{iid}}(N)
    =
    \operatorname{Var}\!\left[
        \frac{1}{N}\sum_{i=1}^{N}\varepsilon_{i,t}
    \right]
    =\frac{1}{N},
\end{equation}
and its local variance elasticity is
\begin{equation}
    \alpha_{\mathrm{iid}}(N)=1.
\end{equation}

Now retain the same weights and marginal variances but impose a common
pairwise correlation $\rho\in(0,1)$. Aggregate variance becomes
\begin{align}
\label{eq:correlated_variance}
    V_{\rho}(N)
    &=
    \frac{1}{N^2}\left[N+\rho N(N-1)\right]
    \nonumber\\
    &=
    \rho+\frac{1-\rho}{N},
\end{align}
with local elasticity
\begin{equation}
\label{eq:correlated_alpha}
    \alpha_{\rho}(N)
    =
    \frac{1-\rho}{\rho N+1-\rho}.
\end{equation}
The independent aggregate continues to diversify at rate $N^{-1}$,
whereas the correlated aggregate approaches the variance floor
$\lim_{N\to\infty}V_{\rho}(N)=\rho$ and
$\lim_{N\to\infty}\alpha_{\rho}(N)=0$.

\begin{figure*}[t]
    \centering
    \includegraphics[
        width=0.98\textwidth
    ]{figures/figure_1_concentration_is_not_scaling.png}
    \caption{
    Concentration and aggregate-variance scaling under equal weights
    and unit marginal shock variance. Panel~(a) compares the common
    concentration path $H(N)=1/N$ with aggregate variance under
    independence and under equicorrelation $\rho=0.15$. The independent
    variance coincides with $H(N)$, whereas the correlated variance
    approaches the floor $\rho$. Panel~(b) reports adjacent-interval
    variance elasticities: concentration and independent variance retain
    elasticity one, while the correlated variance elasticity falls
    toward zero.
    }
    \label{fig:concentration_scaling}
\end{figure*}

On the grid shown in Figure~\ref{fig:concentration_scaling}, the
correlated adjacent-interval variance elasticity falls from approximately
$0.500$ over $N=4$ to $8$ to $0.113$ over $N=32$ to $64$, even though the
concentration elasticity remains exactly one. The two economies have the
same number of units, weights, Herfindahl index, effective number of
firms, and entire concentration path. Yet they have different variance
paths and limiting elasticities. The missing
information is not a refinement of concentration: it is the
scale-dependent variance--covariance structure of the contributions.
This distinction also rules out a common empirical shortcut. Estimating a
power law for $H(N)$ and interpreting its slope as the population exponent
of aggregate variance is valid only after the micro-variance and
covariance terms have been shown to be scale neutral. Concentration and
aggregate-variance scaling are therefore different estimands even when
they coincide under the independent homogeneous benchmark.

This result qualifies rather than overturns the granular mechanism. The
presence of large firms may still make weighted idiosyncratic shocks
quantitatively important. What does not follow is that the population
elasticity of aggregate variance can be recovered from the firm-size
distribution alone. The Herfindahl connection is conditional on stable
homogeneous shock variances, independence, and a coherent population
path. Relaxing those conditions does not merely add noise around the
granular prediction; it changes the estimand that concentration can
identify.

\section{Deterministic analytic verification}
\label{sec:computational_verification}

We verify the accounting chain on the fixed population grid
$N\in\{4,8,16,32,64\}$ using deterministic analytic component paths.
The construction contains no source data, sampling, or stochastic
estimation. For every scenario and population level, the generated
tables record concentration, component values, direct and reconstructed
aggregate variance, adjacent-interval variance elasticities, signed
component contributions, and scale-aware numerical checks. The exercise
therefore tests the implementation of the identities rather than
estimating an economic parameter.

The first attribution example uses two independent positive components,
\begin{equation}
\label{eq:heterogeneous_demo}
    V_{\mathrm{het}}(N)
    =
    N^{-0.9}+0.4N^{-0.3}.
\end{equation}
Their component elasticities differ, so their contributions to the
aggregate variance elasticity change with population scale. The
fast-decaying component supplies most of the initial variance decline,
but its contribution falls as the slower component becomes dominant.
Across the displayed intervals, the reconstructed total variance
elasticity consequently declines from approximately $0.582$ to $0.422$.

The second attribution example groups the equicorrelated counterexample
into total diagonal variance and total doubled covariance,
\begin{align}
\label{eq:grouped_correlated_demo}
    D(N) &= \frac{1}{N}, \\
    B(N) &= \rho\left(1-\frac{1}{N}\right), \\
    V_{\rho}(N) &= D(N)+B(N).
\end{align}
Although $B(N)$ is positive, it rises with population size. Hence
$B(N_1)-B(N_0)>0$ for $N_1>N_0$, and the minus sign in
equation~\eqref{eq:component_alpha_contribution} makes its contribution
to the variance-decay elasticity negative. The covariance channel
therefore offsets part of the diagonal diversification contribution on
every adjacent interval.

\begin{figure*}[t]
    \centering
    \includegraphics[
        width=0.98\textwidth
    ]{figures/figure_2_exact_component_attribution.png}
    \caption{
    Exact component attribution of finite-interval variance
    elasticities. Panel~(a) decomposes the heterogeneous independent
    aggregate into fast- and slow-decaying diagonal contributions.
    Panel~(b) decomposes the correlated equal-weight aggregate into
    diagonal and total doubled-covariance contributions. The positive
    covariance level increases along the population path and therefore
    contributes negatively to variance decay. Open reconstructed-total
    markers overlap the directly computed totals up to floating-point
    precision.
    }
    \label{fig:alpha_verification}
\end{figure*}

The current reference artifact evaluates $15$ population-level variance
identities and $12$ adjacent-interval scaling identities. All pass the
declared scale-aware tolerance contract. Direct and reconstructed
variance coincide exactly in the analytic tables, and the maximum
absolute scaling-reconstruction residual is
$6.7\times10^{-16}$ in the reference run. Such last-digit residuals may
vary across numerical environments; the scientific acceptance criterion
is the stated tolerance, not reproduction of one machine's exact
floating-point remainder. These calculations verify the implementation
and internal consistency of the accounting chain. They do not constitute
empirical evidence about the scaling behavior of any particular economy.

\section{Conclusion}
\label{sec:conclusion}

Concentration locates aggregate exposure; it does not generally identify
how aggregate variance changes as the population expands. The
Herfindahl connection requires a stable homogeneous variance structure,
restricted covariance, and a coherent population path. We provide an
exact finite-interval decomposition that attributes the aggregate
variance elasticity to changes in diagonal and signed covariance
components without logging the components themselves. Dependence can
therefore alter diversification not only by raising aggregate variance
but also through its own population path: a positive covariance
component that rises with $N$ contributes negatively to variance decay.
The result preserves the central granular insight while defining its
boundary. The firm-size distribution can make large firms important,
but it is not a sufficient statistic for population scaling. The
decomposition supplies a transparent accounting contract for empirical
micro-to-macro work, shifting the empirical task from reading an
exponent off the size distribution to measuring how variance and
covariance contributions change along an explicit population path.

% ------------------------------------------------------------------
% Declarations
% ------------------------------------------------------------------

\section*{Funding}

This research did not receive any specific grant from funding agencies
in the public, commercial, or not-for-profit sectors.

\section*{Declaration of competing interest}

The author declares that he has no known competing financial interests
or personal relationships that could have appeared to influence the
work reported in this paper.

\section*{Data and code availability}

No empirical data were used in this study. Source code, deterministic
computational constructions, tests, and generated verification artifacts
will be made available in a public repository.

\section*{
Declaration of generative AI and AI-assisted technologies in the
writing process
}

During the preparation of this work, the author used OpenAI ChatGPT and
Codex to assist with research organization, software implementation,
adversarial checking, and language revision. The author reviewed and
verified the mathematical derivations, computational outputs,
references, and final text, edited the resulting material as needed,
and takes full responsibility for the content of the publication.

% ------------------------------------------------------------------
% References
% ------------------------------------------------------------------

\bibliographystyle{elsarticle-harv}
\bibliography{references}

\end{document}